\documentclass[10pt,a4paper]{amsart}
\usepackage[margin=19mm]{geometry}
\usepackage{amsmath,amssymb,mathtools}
\usepackage[scr=rsfs,cal=euler]{mathalfa}
\usepackage{xfrac}
\usepackage[unicode,hidelinks,bookmarks=false]{hyperref}
\newcommand{\ie}{i.e.\ }
\newcommand{\cf}{cf.\ }
\newcommand{\resp}{respectively\ }
\newcommand{\Subsection}{\subsection}
\providecommand{\nobreakdash}{\nobreak\hbox{-}}
\setlength{\emergencystretch}{2em}
\multlinegap0pt
\usepackage{bm}
\usepackage{bbm}
\usepackage{dsfont}
\usepackage{xspace}
\usepackage{cancel}
\usepackage{mathtools}
\usepackage{amsthm}
\makeatletter
\@ifundefined{teo}{\newtheorem{teo}{Teo}}{}
\@ifundefined{lema}{\newtheorem{lema}[teo]{Lemma}}{}
\@ifundefined{obs}{\newtheorem{obs}[teo]{Remark}}{}
\@ifundefined{propo}{\newtheorem{propo}[teo]{Proposition}}{}
\makeatother
\title[On Matrix Valued Schr\"odinger Operators on the Discrete Rea]{On Matrix Valued Schr\"odinger Operators on the Discrete Real Line: Resolvent Boundary Values, Limiting Absorption Principle, H\"older Regularity and Dispersive Estimates}
\author{Ballesteros Miguel, Franco C\'ordova Gerardo, Gil Jonathan, Ivan Naumkin}
\date{Source: arXiv:2604.01391v1 (CC BY 4.0)}
\begin{document}
\maketitle

\noindent\emph{Source and licence.} On Matrix Valued Schr\"odinger Operators on the Discrete Real Line: Resolvent Boundary Values, Limiting Absorption Principle, H\"older Regularity and Dispersive Estimates. Version arXiv:2604.01391v1 (CC BY 4.0), distributed under the Creative Commons Attribution 4.0 International licence, as recorded in the arXiv licence metadata for this exact version and shown on the abstract page https://arxiv.org/abs/2604.01391v1. Full rights notice: licenses/arxiv-2604.01391-CC-BY-4.0.txt.\par\smallskip

\noindent\textit{Editorial scope.} Counted unit: the Teo whose source label is \texttt{formulaGreenmatrixelements} in the article (source0.tex lines 1694--1708, source label \texttt{formulaGreenmatrixelements}), delivered together with the article's own complete original proof of it (source0.tex lines 1709--1814, one \texttt{proof} environment of 106 lines). No mathematics is written, repaired or re-derived: statement and proof are the article's own text line for line. Required context retained uncounted: eigen (lines 237--239); cauhyproblem (lines 678--680); base (lines 999--1011); RKsr (lines 1404--1418); eq:kernel (lines 1630--1633); aaaa (lines 1636--1655); aaa (lines 1638--1644). Every definition, display and label of those lines is retained unchanged. Omitted from this package and not redistributed: 1--236, 240--677, 681--998, 1012--1403, 1419--1629, 1634--1635, 1645--1693, 1815--2808. Nothing inside a retained excerpt is omitted or altered: every retained body is delivered exactly as the article's lines are written, so no omission receipt of any kind accompanies this package. Citations: the retained excerpts carry no citation command at all, so no bibliography entry of the article is transcribed and no reference is redistributed. Reference adaptation: none was needed and none was made; every reference inside the delivered unit resolves inside it, so no unresolved reference or citation is produced. Printed slips kept literal: no printed word, symbol, hypothesis or number of the counted statement or of its proof was altered. Layout and class: the article's own class is \texttt{article} with packages including inputenc, babel, amsmath, amsthm, amsfonts, amssymb, graphicx, hyperref, geometry, color, appendix, authblk; the wrapper is the lane's pinned \texttt{amsart} editorial wrapper at 10pt on A4, which restates the theorem-like environments the retained text needs and adds the article's own macro definitions that the retained text needs (none), with extra wrapper packages (bm, bbm, dsfont, xspace, cancel, mathtools). The delivered numbering is the unit's own. Assets: no retained span contains a figure environment, a graphics-inclusion command, a PDF-page inclusion, a table or a TikZ picture; no image file is copied into this package, the package declares no asset dependency and no image file is redistributed with it. Licence: version arXiv:2604.01391v1 of this article is distributed under the Creative Commons Attribution 4.0 International licence, as recorded in the arXiv licence metadata for this exact version and shown on the abstract page https://arxiv.org/abs/2604.01391v1. A full rights notice is delivered at licenses/arxiv-2604.01391-CC-BY-4.0.txt. Characters: every retained excerpt is pure ASCII, so the delivered engine, XeLaTeX with its default Latin Modern text font, reports no missing character.
\begin{align}\label{eigen}
 \tau u = E u, \qquad \tau_0 u = E u,
 \end{align}

 \begin{teo}\label{cauhyproblem}
   Let \(A_{0}, B_{0} \in \mathcal{M}_{L}\) and \(n_{0} \in \mathbb{Z}\). There exists a unique solution \(u\) of the eigenvalue equation \eqref{eigen} satisfying \(u(n_{0}) = A_{0}\) and \(u(n_{0}+1) = B_{0}\).
 \end{teo}

\begin{propo}\label{base}
For $z\in \overline{\mathbb{D}}\setminus \{-1, 0, 1\},$ and $E=J(z)$, any solution of \eqref{eigen} can be uniquely written as 
\begin{align}\label{yh}
s^{z} = u^{z}_+ A_+ + v^{z^{- 1}}_+ B_+, \hspace{1cm} 
s^{z} = u^{z^{ - 1}}_- A_- + v^{z}_- B_-
\end{align} 
where $A_\pm,B_\pm \in \mathcal{M}_{L}$. 
%Moreover, any solution of \eqref{eigen} can also be uniquely written as
%\begin{align}\label{yh2}
%s^{z} = u^{z^{ - 1}}_- A_- + v^{z}_- B_-,
%\end{align}
%where $A_-,B_-\in \mathcal{M}_{L}$.
\end{propo}

\begin{obs}\label{RKsr}
For every $s$ and $r$, the sequences $([\mathcal{K}]_{s,x})_{x \in \mathbb{Z}}$ and $([\mathcal{K}]_{x,r})_{x \in \mathbb{Z}}$ are square-summable. Indeed, for any $A \in \mathcal{M}_L$, define $u_{r,A} := \pi^{*}_{r}(A) \in \ell^{2}(\mathbb{Z},\mathcal{M}_{L})$. Then
\begin{align}\label{KER}
\|([\mathcal{K}]_{x,r}A)_{x \in \mathbb{Z}}\|_{\ell^{2}} = \|\mathcal{K}u_{r,A}\|_{\ell^{2}} \leq \|\mathcal{K}\| \|A\|_{\mathcal{M}_L}.
\end{align}
This implies that $([\mathcal{K}]_{x,r})_{x \in \mathbb{Z}}$ is square-summable. The square-summability of $([\mathcal{K}]_{s,x})_{x \in \mathbb{Z}}$ follows by noting that
$ 
[\mathcal{K}]_{s,x} = \left([\mathcal{K}^*]_{x,s}\right)^*.
$ 

Consequently, for every $u \in \ell^{2}(\mathbb{Z},\mathcal{M}_{L})$, we have $\sum_{r\in \mathbb{Z}} \|[\mathcal{K}]_{s,r}u(r)\|_{\mathcal{M}_L} < \infty$ and
\begin{align}
(\mathcal{K}u)(s) = \sum_{r\in \mathbb{Z}} [\mathcal{K}]_{s,r}u(r), \quad u \in \ell^{2}(\mathbb{Z},\mathcal{M}_{L}).
\end{align}
\end{obs}

\begin{equation}
\label{eq:kernel}
[R_{H}(E)]_{s+1,r} + [R_{H}(E)]_{s-1,r} + (V(s) - E) [R_{H}(E)]_{s,r} = \pi_s \pi_r^{*}.
\end{equation}

\begin{lema}\label{aaaa}
    For every $E \in \mathbb{C}\setminus \sigma(H)$, where $E = J(z)$, the matrix
    \begin{align}\label{aaa}
        \Phi(u_+^z, u_-^{z^{-1}})(r) = 
        \begin{pmatrix}
            u_+^z(r+1) & u_-^{z^{-1}}(r+1) \\
            u_+^z(r) & u_-^{z^{-1}}(r)
        \end{pmatrix}
    \end{align}
    is invertible, and its inverse is given by:
    \begin{align}
        \Phi(u_+^z, u_-^{z^{-1}})(r)^{-1} = 
       i \begin{pmatrix}
            -W(u_-^{\overline{z}^{-1}}, u_+^z)^{-1} u_-^{\overline{z}^{-1}}(r)^* 
            & W(u_-^{\overline{z}^{-1}}, u_+^z)^{-1} u_-^{\overline{z}^{-1}}(r+1)^* \\
            -W(u_+^{\overline{z}}, u_-^{z^{-1}})^{-1} u_+^{\overline{z}}(r)^* 
            & W(u_+^{\overline{z}}, u_-^{z^{-1}})^{-1} u_+^{\overline{z}}(r+1)^*
        \end{pmatrix}.
    \end{align}
\end{lema}

\begin{teo}\label{formulaGreenmatrixelements}
    For  $E\in \mathbb{C}\setminus \sigma (H)$, the kernel of $R_{H}$ is given by
    \begin{equation}
    [R_{H}(E)]_{s,r}=
    \label{eq:Green}
 \left\{ \begin{array}{lcc}
               -iu_{+}^{z}(s) W(u_{-}^{\overline{z}^{-1}},u_{+}^{z},)^{-1} u_{-}^{\overline{z}^{-1}}(r)^{*} &   s\geq r \\
            iu_{-}^{z^{-1}}(s) W(u_{+}^{\overline{z}}, u_{-}^{z^{-1}})^{-1} u_{+}^{\overline{z}}(r)^{*} &  s< r
             \end{array}
   \right.
   \end{equation}
   for $J(z)=E$ where $z\in \mathbb{D}\setminus \{0\}$. 
   
   The matrix elements, $[R_{H}(E)]_{s,r}$, are called Green functions.
\end{teo}

\begin{proof}
   Regarding equation \eqref{eq:kernel}, the following holds for all $C\in \mathcal{M}_L$:
\begin{equation}
     \label{eq:kernel2}
     [R_{H}(E)]_{ s+1,r}C + [R_{H}(E)]_{ s-1,r}C + (V(s) - E)[R_{H}(E)]_{s,r}C = \delta_{r,s}  C.
\end{equation}
In particular, for \(s = r\), we obtain:
\begin{equation}\label{deq}
     [R_{H}(E)]_{r+1,r}C + [R_{H}(E)]_{r-1,r}C + (V(r) - E)[R_{H}(E)]_{r,r}C = C.
\end{equation}

On the other hand, for a fixed \(r \in \mathbb{Z}\), the sequence \( s \mapsto [R_{H}(E)]_{s,r}C\), for $s>r$, satisfies the equation
\begin{align}
[R_{H}(E)]_{ s+1,r}C + [R_{H}(E)]_{ s-1,r}C + (V(s) - E)[R_{H}(E)]_{s,r}C=0 . 
\end{align}
The solution to the Cauchy problem 
\begin{align}
    \tau w = E w ,   \hspace{1cm}  w(r+1) = [R_{H}(E)]_{r+1,r}C,   \hspace{.3cm}
     w(r+2) = [R_{H}(E)]_{r+2,r}C
     \end{align}
is unique and it is given in Theorem \ref{cauhyproblem}. Then, $w$ is a unique extension of 
   \( \Big (  [R_{H}(E)]_{s,r}C \Big )_ {s>r}\),  to $ \mathbb{Z} $. 
 Recall that by Proposition \ref{base}, there exist matrices \(A^{z}, B^{z} \in \mathcal{M}_{L}\) such that:
\begin{align}
    w(s) = u_{+}^{z}(s) A^{z} + v_{+}^{z^{-1}}(s) D^{z}.
\end{align}
Since the sequence \( s \mapsto [R_{H}(E)]_{s,r}C\) is square summable (see Remark \ref{RKsr})  and \(v_{+}^{z^{-1}}(s) = z^{-s}(I + o(1))\) as \(n \to \infty\), we have that:
 \begin{align}\label{nj}
  w(s) =u_{+}^{z}(s) A^{z}, \:  \: \forall s,  \hspace{.5cm}   [R_{H}(E)]_{s,r}C =u_{+}^{z}(s) A^{z},\,\,\, s\geq r+1  
 \end{align}
($A^z$ and $D^z$ depend on $r$, but for simplicity in the proof, we will omit it).

Substituting \(s = r+1\) in \eqref{eq:kernel2}, we obtain:
\begin{align}\label{g}
     [R_{H}(E)]_{r+2,r} + [R_{H}(E)]_{r,r} + (V(r+1) - E)[R_{H}(E)]_{ r+1,r} = 0
\end{align}
  and this equation is equivalent to 
  \begin{align}\label{g1}
      [R_{H}(E)]_{r,r} =-  [R_{H}(E)]_{r+2,r}-(V(r+1) - E)[R_{H}(E)]_{r+1,r}.
\end{align}
Now we use the fact that $w$ satisfies the equation  \eqref{nj}, that  $\tau w=Ew$, that $  w $ is an extension   of  \( \Big (  [R_{H}(E)]_{s,r}C \Big )_ {s>r}\) and  \eqref{g1} to obtain 
    \begin{align}\label{nj2}
u_{+}^{z}(r) A^{z} = w(r) = - w(r+2) -(V(r+1) - E ) w(r+1) =      [R_{H}(E)]_{r,r}C .
\end{align}
From equations \eqref{nj} and \eqref{nj2}, we conclude:
\begin{align} \label{nj3}
     [R_{H}(E)]_{s,r}C = u_{+}^{z}(s) A^{z}, \quad \text{for } s \geq r.
\end{align}
A similar analysis yields there there is a matrix $B^z$ such that 
\begin{align} \label{nj4}
     [R_{H}(E)]_{s,r}C= u_{-}^{z^{-1}}(s) B^{z}, \quad \text{for } s \leq r.
\end{align}

Now, substituting \eqref{nj3} and \eqref{nj4} into \eqref{deq} for \(s = r\), we get:
\begin{align}\label{nj5}
     u_{+}^{z}(r+1) A^{z} + u_{-}^{z^{-1}}(r-1) B^{z} + (V(r) - E) u_{-}^{z^{-1}}(r) B^{z} = C.
\end{align}
Since \(u_{-}^{z^{-1}}\) satisfies \(\tau u = Eu \), equation \eqref{nj5} simplifies to:
\begin{align}\label{s1}
     u_{+}^{z}(r+1) A^{z} - u_{-}^{z^{-1}}(r+1) B^{z} = C.
\end{align}
Additionally, from \eqref{nj3} and \eqref{nj4} evaluated at \(s = r\), we have:
\begin{align}\label{s2}
     u_{+}^{z}(r) A^{z} = u_{-}^{z^{-1}}(r) B^{z}.
\end{align}
Finally, from \eqref{s1} and \eqref{s2}, we obtain the following system of equations:
\begin{align}
     u_{+}^{z}(r+1) A^{z} - u_{-}^{z^{-1}}(r+1) B^{z} &= C, \\
     u_{+}^{z}(r) A^{z} - u_{-}^{z^{-1}}(r) B^{z} &= 0 .
\end{align}
This can be writen in the next form: 
\begin{align}
    \Phi ( u_{+}^{z}, u_{-}^{z^{-1}})(r)
    \begin{pmatrix}
        A^{z}\\
        -B^{z}
         \end{pmatrix}=
          \begin{pmatrix}
        C\\
        0
         \end{pmatrix},
\end{align}
where we use \eqref{aaa}.  Now we recall  Lemma \ref{aaaa} to obtain
\begin{align}
    \begin{pmatrix}
        A^{z}\\
       - B^{z}
         \end{pmatrix}= i\begin{pmatrix}
           -W(u_- ^{\overline{z}^{-1}}, u_+^{z})^{-1}u_- ^{\overline{z}^{-1}}(r)^{*} & W(u_- ^{\overline{z}^{-1}}, u_+^{z})^{-1}u_- ^{\overline{z}^{-1}}(r+1)^{*} \\
           -W(u_+ ^{\overline{z}}, u_-^{z^{-1}})^{-1}u_+ ^{\overline{z}}(r)^{*} & W(u_+ ^{\overline{z}}, u_-^{z^{-1}})^{-1}u_+ ^{\overline{z}}(r+1)^{*}
        \end{pmatrix} \begin{pmatrix}
        C\\
        0
         \end{pmatrix}.
\end{align}

Finally, we use  \eqref{nj3} and \eqref{nj4} to get
 \begin{equation}
    [R_{H}(E)]_{s,r}C=
 \left\{ \begin{array}{lcc}
             -iu_{+}^{z}(s) W(u_{-}^{\overline{z}^{-1}},u_{+}^{z},)^{-1} u_{-}^{\overline{z}^{-1}}(r)^{*}C &   s\geq r \\
            iu_{-}^{z^{-1}}(s) W(u_{+}^{\overline{z}}, u_{-}^{z^{-1}})^{-1} u_{+}^{\overline{z}}(r)^{*}C &  s< r
             \end{array}
   \right.
   \end{equation}  for all $C\in \mathcal{M}_L.$
\end{proof}

\end{document}
